\section{Exact return means and stationary unfinished flights}
\label{sec:exact-return-means}
The enlarged section gives explicit normalization constants without any
local-limit assumption. We use ergodicity of the finite-horizon
periodic dispersing collision map, as supplied by the standard billiard
framework \cite{Young}. No spectral conclusion for the new first-return
map is imported from that fact.

\begin{proposition}[Exact induced means]\label{prop:exact-induced-means}
For the actual return records to $Y_R^*$,
\begin{equation}\label{eq:induced-mean-vector}
 \int(\kappa_{R,1}^*,\kappa_{R,2}^*,r_R^*,\tau_R^*)\dd\nu_R^*
 =\left(0,0,\frac1{c_*},\frac{\bar\tau_R}{c_*}\right),
 \quad
 \bar\tau_R=\frac{\sqrt3}{4R}-\frac{\pi R}{2}.
\end{equation}
The mean roof satisfies
\begin{equation}\label{eq:mean-roof-modulus}
 \left|\frac{\mathrm d}{\mathrm dR}\int\tau_R^*\dd\nu_R^*\right|
 <\frac{15}{4c_*}\qquad(R\in I).
\end{equation}
The equilibrium physical collision rate remains $1/\bar\tau_R$.
\end{proposition}
\begin{proof}
For any integrable physical collision observable $f$, the disjoint
return-tower levels partition the collision section up to null sets.
Invariance, first for nonnegative functions and then for positive and
negative parts, gives the Kac identity
\[
 \int_{Y_R^*}\sum_{j=0}^{r_R^*(x)-1} f(T_R^j x)\dd\nu(x)
                         =\int f\dd\nu.
\]
Taking $f=1$ and dividing by $c_*$ gives the mean return count.
Taking $f=\tau_R$ gives the mean induced roof. The physical
suspension normalization proved earlier gives the displayed value
of $\bar\tau_R$. The physical displacement is integrable by the
finite-horizon bound. Spatial inversion of the entire lattice maps
its value to its negative and preserves $\nu$, so its mean is zero.
This proves \eqref{eq:induced-mean-vector}.

The section volume $c_*$ is independent of $R$. Differentiating the
explicit expression for $\bar\tau_R$ yields
$-\sqrt3/(4R^2)-\pi/2$. The bounds $\sqrt3<7/4$,
$\pi<22/7$ and $R\ge9/20$ show its absolute value is less than
$15/4$. Finally the ratio of mean physical collisions per return to
mean physical time per return is $1/\bar\tau_R$.
\end{proof}

\begin{proposition}[The complete induced suspension]\label{prop:induced-suspension}
The equilibrium flow is represented by the suspension over $F_R^*$
with roof $\tau_R^*$ and probability
\begin{equation}\label{eq:induced-suspension}
 \frac{\nu_R^*(\mathrm dx)\,\mathrm ds}{\int\tau_R^*\dd\nu_R^*},
               \qquad 0\le s<\tau_R^*(x).
\end{equation}
If $A_R$ and $B_R$ denote respectively the age since the last visit
and the time to the next visit to $Y_R^*$ under this stationary law,
then, for $q\ge0$,
\begin{equation}\label{eq:residual-tails}
 \Pr(A_R>q)=\Pr(B_R>q)
 =\frac{\int(\tau_R^*-q)_+\dd\nu_R^*}
              {\int\tau_R^*\dd\nu_R^*}.
\end{equation}
At either one fixed endpoint of a time interval of length $t$,
the incomplete induced flight, its physical collision count and its
lattice displacement divided by $\sqrt t$ converge to zero in
probability as $t\to\infty$, for each fixed $R$.
\end{proposition}
\begin{proof}
In the physical suspension over $T_R$, group the consecutive
single-flight intervals before each first return to $Y_R^*$.
The tower identity and the normalization
$c_*\int\tau_R^*\dd\nu_R^*=\bar\tau_R$ convert the original
suspension measure exactly to \eqref{eq:induced-suspension}.
This rearranges the actual dependent path; it introduces no new
launch law and no independent flight variables. The marked formula
of Proposition~\ref{prop:marked} can consequently be grouped into
these complete return records while keeping both endpoint pieces.

For a given roof $\tau$, the set of $s\in[0,\tau)$ with age $s>q$
has length $(\tau-q)_+$. The set with remaining time $\tau-s>q$
has the same length. Integration proves \eqref{eq:residual-tails}.
Its right side tends to zero by the integrability in
Proposition~\ref{prop:exact-induced-means}. Each endpoint roof is
finite almost surely under the length-biased stationary law;
stationarity gives the same marginal at the other fixed endpoint.
Thus its ratio to $\sqrt t$ tends to zero in probability.

Every physical flight has length at least $1-2R\ge3/50$. The
number of collisions in an incomplete interval of length $s$ is
at most $1+s/(1-2R)$. Its displacement between fundamental cells
is bounded by a constant plus a fixed multiple of $s$, since the
speed is one and the lattice cell has bounded diameter. These
bounds prove the same endpoint statement for the two marked rewards.
\end{proof}

A fixed-endpoint assertion is not a bound on the maximum residual over
all times in $[0,t]$. Nor does pointwise integrability in $R$ imply a
uniform residual estimate over $I$. Such uniform or process-level
claims require uniform integrability or a quantitative tail for the
induced roof. The functional limit assumptions in
Proposition~\ref{prop:clock} therefore remain explicit. The exact
mean and the modulus \eqref{eq:mean-roof-modulus}, by contrast, are
already available for geometric error propagation. They do not
require recovery of the entire preparation distribution in A2-GEOM.
